% p.22 — الباب الثاني العشر في باع يبيع (اجوف يائي)
\mspage{22}
\bab{\com{الباب الثاني العشر}{كذا في الأصل؛ ولعل الصواب: الباب الثاني عشر} في باع يبيع (اجوف يائي)}
% Layout: six tall columns (right-to-left), each a stack of boxed blocks of unequal height.
% Column 1 is the summary column: the باب name (ضَرَبَ) in its own box, then the ten
% underlying (pre-iʿlāl) forms بَيَعَ … لَا تَبْيِعْ, then the ten actual forms بَاعَ … لَا تَبِعْ.
% In columns 2–6 every sub-heading box holds the heading and, beneath it, the underlying form
% (e.g. «ماضي معروف / بَيَعَ»); the regular paradigm follows outside the box.
% Row alignment is approximate (all columns run to 21 lines here); \cline marks the printed box
% boundaries inside each column.
\begin{sarf}{|c|c|c|c|c|c|}
ضَرَبَ & \thead{ماضي معروف} & تَبِيْعُ & بِيْعَ & يُبَعْنَ & بِيْعُوْا\\\cline{1-1}
بَيَعَ & بَيَعَ & تَبِيْعَانِ & بِيْعَا & تُبَاعُ & بِيْعِيْ\\\cline{2-2}
يَبْيِعُ & بَاعَ & يَبِعْنَ & بِيْعُوْا & تُبَاعَانِ & بِيْعَا\\
بَيْعًا & بَاعَا & تَبِيْعُ & بِيْعَتْ & تُبَاعُوْنَ & بِعْنَ\\\cline{6-6}
بَايِعٌ & بَاعُوْا & تَبِيْعَانِ & بِيْعَتَا & تُبَاعِيْنَ & \thead{نهي حاضر}\\
بُيِعَ & بَاعَتْ & تَبِيْعُوْنَ & بِعْنَ & تُبَاعَانِ & لَا تَبْيِعْ\\\cline{6-6}
يُبْيَعُ & بَاعَتَا & تَبِيْعِيْنَ & بِعْتَ & تُبَعْنَ & لَا تَبِعْ\\
بَيْعًا & بِعْنَ & تَبِيْعَانِ & بِعْتُمَا & أُبَاعُ & لَا تَبِيْعَا\\
مَبْيُوْعٌ & بِعْتَ & تَبِعْنَ & بِعْتُمْ & نُبَاعُ & لَا تَبِيْعُوْا\\\cline{5-5}
اِبْيِعْ & بِعْتُمَا & أَبِيْعُ & بِعْتِ & \thead{اسم مفعول} & لَا تَبِيْعِيْ\\
لَا تَبْيِعْ & بِعْتُمْ & نَبِيْعُ & بِعْتُمَا & مَبْيُوْعٌ & لَا تَبِيْعَا\\\cline{1-1}\cline{3-3}\cline{5-5}
بَاعَ & بِعْتِ & \thead{اسم فاعل} & بِعْتُنَّ & مَبِيْعٌ & لَا تَبِعْنَ\\\cline{6-6}
يَبِيْعُ & بِعْتُمَا & بَايِعٌ & بِعْتُ & مَبِيْعَانِ & \thead{مضارع بلم}\\\cline{3-3}
بَيْعًا & بِعْتُنَّ & بَائِعٌ & بِعْنَا & مَبِيْعُوْنَ & لَمْ يَبْيِعْ\\\cline{4-4}\cline{6-6}
بَائِعٌ & بِعْتُ & بَائِعَانِ & \thead{مضارع مجهول} & مَبِيْعَةٌ & لَمْ يَبِعْ\\
بِيْعَ & بِعْنَا & بَائِعُوْنَ & يُبْيَعُ & مَبِيْعَتَانِ & لَمْ يَبِيْعَا\\\cline{2-2}\cline{4-4}
يُبَاعُ & \thead{مضارع معروف} & بَائِعَةٌ & يُبَاعُ & مَبِيْعَاتٌ & لَمْ يَبِيْعُوْا\\\cline{5-5}
بَيْعًا & يَبْيِعُ & بَائِعَتَانِ & يُبَاعَانِ & \thead{أمر حاضر} & لَمْ تَبِعْ\\\cline{2-2}
مَبِيْعٌ & يَبِيْعُ & بَائِعَاتٌ & يُبَاعُوْنَ & اِبْيِعْ & لَمْ تَبِيْعَا\\\cline{3-3}\cline{5-5}
بِعْ & يَبِيْعَانِ & \thead{ماضي مجهول} & تُبَاعُ & بِعْ & لَمْ يَبِعْنَ\\
لَا تَبِعْ & يَبِيْعُوْنَ & بُيِعَ & تُبَاعَانِ & بِيْعَا & لَمْ تَبِعْ\\
\end{sarf}
